\section{可三角化自同态}

记号约定:当$V$是有限维$K$-向量空间,$\mathcal{T}$的所有特征值都属于$K$时,记$\mathscr{E}$是$\mathcal{T}$在$E$上的特征值构成的集合.

\begin{definition}{广义特征向量,广义特征子空间}{}
	设$V$是有限维$K$-向量空间,$\lambda\in K,n\geq 1$.
	\begin{enumerate}
		\item 我们称$\GEig_{K,n}\left(V,\mathcal{T},\lambda\right)\coloneq\ker\left(\left(\lambda\id_{V}-\mathcal{T}\right)^{n}\right)$是$\lambda$的$n$阶广义特征子空间.
		\item 当$\GEig_{K,n}\left(V,\mathcal{T},\lambda\right)\neq\left\{0\right\}$时,称$\GEig_{K,n}\left(V,\mathcal{T},\lambda\right)\setminus\left\{0\right\}$的元素为$\lambda$的$n$阶广义特征向量.
		\item 我们称$\GEig_{K}\left(V,\mathcal{T},\lambda\right)\coloneq\bigcup\limits_{k=1}^{\infty}\GEig_{K,k}\left(V,\mathcal{T},\lambda\right)$是$\lambda$的广义特征子空间.
		\item 我们称$\gm\left(\lambda\right)\coloneq\dim_{K}\GEig_{K}\left(V,\mathcal{T},\lambda\right)$为$\lambda$的代数重数.
	\end{enumerate}
\end{definition}

\begin{definition}{Jordan块}{}
	对每个$\alpha\in K$和$n\geq 1$,称$J_{n}\left(\alpha\right)\coloneq\alpha\boldsymbol{I}_{n,n}+\sum\limits_{i=1}^{n-1}\boldsymbol{E}_{i+1,i}$是$n$阶Jordan块,称$n$是$J_{n}\left(\alpha\right)$的阶数.
\end{definition}

\begin{proposition}{向量空间的广义特征子空间分解}{general-eigensubspaces-decomposition-of-vector-spaces}
	设$V$是有限维$K$-向量空间,$\mathcal{T}$的所有特征值都属于$K$.
	\begin{enumerate}
		\item $V=\bigoplus\limits_{\lambda\in\mathscr{E}}\GEig_{K}\left(V,\mathcal{T},\lambda\right)$.
		\item 对每个$\lambda\in\mathscr{E}$,记$p_{\lambda}:V\to\GEig_{K}\left(V,\mathcal{T},\lambda\right)$是典范映射.那么$\forall\lambda\in\mathscr{E}\exists s_{\lambda}\in K[X]\forall x\in V\left(p_{\lambda}\left(x\right)=s_{\lambda}\left(\mathcal{T}\right)x\right)$.
	\end{enumerate}
\end{proposition}

\begin{proposition}{自同态的Jordan块的构造}{}
	设$V$是有限维$K$-向量空间,$\mathcal{T}$的所有特征值都属于$K$,则存在$r\geq 1$和$V$的子空间族$\left(V_{i}\right)_{i=1}^{r}$满足如下条件:
	\begin{enumerate}
		\item $V=\bigoplus\limits_{i=1}^{r}V_{i}$,并且对每个$1\leq i\leq r$,$E_{i}$是$\mathcal{T}$-不变子空间和$\mathcal{T}$-循环子空间.
		\item 对每个$1\leq i\leq r$,记$\mathcal{T}_{i}:E_{i}\to E_{i},x\mapsto\mathcal{T}\left(x\right)$.那么,对每个$1\leq i\leq r$,存在$\alpha_{i}\in\mathscr{E}$和$m\left(i\right)\geq 1$使得
		\begin{align*}
			\Min_{\mathcal{T}_{i},V_{i}}=\left(X-\lambda_{i}\right)^{m_{i}}.
		\end{align*}
		对每个$1\leq i\leq r$,记$p_{i}=\Min_{\mathcal{T}_{i},V_{i}},\mathfrak{p}_{m\left(i\right)}=\langle p_{i}\rangle$.
		\item 对每个$1\leq i\leq r$,把$V_{i}$视为$K[X]$-模后得到的模记作$\tilde{V}_{i}$,$\tau_{i}:K[X]\to K[X]/\mathfrak{p}_{m\left(i\right)}$是典范同构.则对每个$1\leq i\leq r$,
		\begin{align*}
			\tilde{\mathcal{B}}_{i}\coloneq((X-\lambda_{i})^{k}+\mathfrak{p}_{m\left(i\right)})_{i=0}^{m\left(i\right)-1},
			\mathcal{B}_{i}\coloneq(\tau_{i}^{-1}((X-\lambda_{i})^{k}+\mathfrak{p}_{m\left(i\right)}))_{i=0}^{m\left(i\right)-1}
		\end{align*}
		分别是$\tilde{V}_{i},V_{i}$的基.
		\item 对每个$1\leq i\leq r$有$\left[\mathcal{T}_{i}\right]_{\mathcal{B}_{i}}^{\mathcal{B}_{i}}=J_{m_{i}}\left(\lambda_{i}\right)$.
	\end{enumerate}
\end{proposition}

\begin{proposition}{有限维向量空间上线性自同态可三角化的等价条件}{equivalent-condition-of-triangularisable-endomorphism-over-finite-dimensional-vector-spaces}
	设$V$是$n$维$K$-向量空间.下列陈述等价:
	\begin{enumerate}
		\item 存在$K$的代数闭扩域$L$,使得$\mathcal{T}$在$L$中的所有特征值都属于$K$.
		\item 存在$V$的基$\mathcal{B}\coloneq\left(e_{i}\right)_{i=1}^{n}$使得$\left[\mathcal{T}\right]_{\mathcal{B}}^{\mathcal{B}}\in \mathfrak{b}_{n}\left(K\right)\cup\mathfrak{b}_{n}^{-}\left(K\right)$.
		\item 存在$V$的基$\mathcal{C}\coloneq\left(\check{e}_{i}\right)_{i=1}^{n}$,$m\geq 1$和$K\times\N_{\geq 1}$的序列$\left(\left(\lambda_{i},k_{i}\right)\right)_{i=1}^{m}$使得
		\begin{align*}
			\left[\mathcal{T}\right]_{\mathcal{C}}^{\mathcal{C}}=\diag\left(J_{k_{1}}\left(\lambda_{1}\right),\ldots,J_{k_{m}}\left(\lambda_{m}\right)\right),k_{1}+\ldots+k_{m}=n.
		\end{align*}
	\end{enumerate}
\end{proposition}

\begin{definition}{可三角化自同态}{}
	沿用命题(\ref{prop:equivalent-condition-of-triangularisable-endomorphism-over-finite-dimensional-vector-spaces})的设定.其中的三个陈述之一为真时,称$\mathcal{T}$可三角化.
\end{definition}

\begin{remark}{}{}
	显然,当$K$是代数闭域时,$\mathcal{T}$一定可三角化.
\end{remark}

\begin{remark}{}{}
	设$m\geq 1$,$\left(\left(\lambda_{i},k_{i}\right)\right)_{i=1}^{m}$是$K\times\N_{\geq 1}$中的序列,$\boldsymbol{A}\in\mathbf{M}_{n}\left(K\right)$且和$K$上的矩阵
	\begin{align*}
		\boldsymbol{J}\coloneq\diag\left(J_{k_{1}}\left(\lambda_{1}\right),\ldots,J_{k_{m}}\left(\lambda_{m}\right)\right)
	\end{align*}
	相似.那么,对$\boldsymbol{A}$在$K$中的每个特征值$\alpha$和每个$m\geq 1$,矩阵$\boldsymbol{J}$中的$m$阶Jordan块$J_{m}\left(\alpha\right)$的个数是唯一的.换言之:
	\begin{center}
		在精确到对角元的排列的意义下,矩阵$\boldsymbol{J}$是唯一的.
	\end{center}
	
	我们也可以采用之前提到的方法来计算$\boldsymbol{A}$的相似不变量:
	\begin{itemize}
		\item 对应于同一个特征值的所有初等因子放在同一行.在每一行中,把这些多项式按照幂次(降序)从左向右排列.
		\item 在每行的右侧用$1$进行填补,使得所有行的长度一致.当补上$1$后的项数等于矩阵的阶数(或初等因子个数最多的那一行的长度)时,停止补$1$.
		\item 把同一列的所有元素相乘,从左到右依次读取每一列的乘积,如此得到的多项式就是矩阵$\boldsymbol{A}$的相似不变量.
	\end{itemize}
	例如,对于矩阵$
	\begin{bmatrix}
		2&0&0\\
		0&3&0\\
		0&1&3
	\end{bmatrix}$,写下
	\begin{gather*}
		\begin{array}{rrr}
			\left(X-2\right),&1,&1\\
			\left(X-3\right)^{2},&1,&1
		\end{array}
	\end{gather*}	
	那么相似不变量就是$1,1,(X-2)(X-3)^{2}$.
\end{remark}

\begin{proposition}{矩阵的极小多项式的计算}{}
	设$\boldsymbol{A}\in\mathbf{M}_{n}\left(K\right),n\geq 1$.如果存在$m\geq 1$和$K\times\N_{\geq 1}$的序列$\left(\left(\lambda_{i},k_{i}\right)\right)_{i=1}^{m}$,使得
	\begin{align*}
		\boldsymbol{A}\text{和}\diag\left(J_{k_{1}}\left(\lambda_{1}\right),\ldots,J_{k_{m}}\left(\lambda_{m}\right)\right)\text{相似},k_{1}+\ldots+k_{m}=n.
	\end{align*}
	那么$\Min_{\boldsymbol{A}}$是$((X-\lambda_{i})^{k_{i}})_{i=1}^{m}$的最小公倍数,$\Char_{\boldsymbol{A}}=\prod\limits_{i=1}^{m}\left(X-\lambda_{i}\right)^{m_{i}}$.
\end{proposition}

\begin{corollary}{代数重数的几何含义}{}
	沿用命题(\ref{prop:general-eigensubspaces-decomposition-of-vector-spaces})的设定,则对每个$\lambda\in\mathscr{E}$,它作为$\Char_{\mathcal{T}}$的根的重数和$\gm\left(\lambda\right)$相等.
\end{corollary}
